\documentclass[10pt,a4paper]{amsart}
\usepackage[margin=19mm]{geometry}
\usepackage{amsmath,amssymb,mathtools}
\usepackage[scr=rsfs,cal=euler]{mathalfa}
\usepackage{xfrac}
\usepackage[unicode,hidelinks,bookmarks=false]{hyperref}
\newcommand{\ie}{i.e.\ }
\newcommand{\cf}{cf.\ }
\newcommand{\resp}{respectively\ }
\newcommand{\Subsection}{\subsection}
\providecommand{\nobreakdash}{\nobreak\hbox{-}}
\setlength{\emergencystretch}{2em}
\multlinegap0pt
\usepackage{amssymb}
\usepackage{tikz}
\usepackage{bm}
\usepackage{mathtools}
\usepackage{enumerate}
\usepackage[shortlabels]{enumitem}
\newtheorem{lemma}{Lemma}
\newcommand{\LAG}{\textup{DUAL}}
\newcommand{\OPT}{\textup{OPT}}
\newcommand{\bv}[0]{\mathbf{b}}
\DeclareMathOperator{\conv}{\operatorname{conv}}
\newcommand{\cv}[0]{\mathbf{c}}
\newcommand{\iprod}[2]{\left\langle {#1}, {#2} \right\rangle}
\newcommand{\xv}[0]{\mathbf{x}}
\title{Probabilistic analysis of dual decomposition on two-stage stochastic integer programs}
\author{Santanu S. Dey, Marco Molinaro, Jingye Xu}
\date{Source: arXiv:2604.23383v1 (CC BY 4.0)}
\begin{document}
\maketitle

\noindent\emph{Source and licence.} Probabilistic analysis of dual decomposition on two-stage stochastic integer programs. Version arXiv:2604.23383v1 (CC BY 4.0), distributed by arXiv under the Creative Commons Attribution 4.0 International licence, http://creativecommons.org/licenses/by/4.0/, as recorded in the arXiv licence metadata for this exact version and shown on the abstract page https://arxiv.org/abs/2604.23383v1. Full licence notice: \texttt{licenses/arxiv-2604.23383-CC-BY-4.0.txt}.\par\smallskip

\noindent\textit{Editorial scope.} Counted unit: the article's lemma lemma:bestBound (source0.tex lines 889--891). The article's complete original proof of it is delivered with it (source0.tex lines 893--963, one \texttt{proof} environment of 71 lines). No mathematics is written, repaired or re-derived: the statement and the proof are the article's own lines, in the article's own notation, and no retained line is substituted or re-flowed.  Retained uncounted as the source-local context the proof needs: lines 736--742.  Nothing was omitted from any retained span, so the package carries no omission record.  No retained line contains a citation, so the wrapper carries no bibliography and no bibliography entry.  No retained line contains a figure environment, an image-inclusion command or drawing code, so no image is delivered and the package declares no asset dependency.  Editorial formatting only, in addition: an AMS \texttt{amsart} preamble that restates the article's own declarations and macros used by the retained lines, namely the environments \texttt{lemma} on separate counters, together with the standard packages those lines need.  The retained environments print in this document as Lemma 1 in order of appearance, whereas the article prints the same items with the numbering of its own full document.  The complete native source of the article as distributed in the arXiv e-print for this exact version is bundled unchanged as \texttt{sources/199161/source0.tex}.
\begin{equation}
    \label{eq:lagrelax}
    \begin{aligned}
    \LAG := \max_{\xv} \ & s \cdot \iprod{\cv^{(0)}}{\xv^{(0)}} + \sum_{i \in [s]} \iprod{\cv^{(i)}}{\xv^{(i)}} \\
    \text{s.t.}\ & (\xv^{(0)}, \xv^{(i)}) \in \conv(\mathcal{X}^{(i)}), \quad \forall i \in [s],
    \end{aligned}
\end{equation}

\begin{lemma} \label{lemma:bestBound}
    Consider any internal node in the Branch-and-Price tree, and let $\bar{\xv}^{(0)},\bar{\xv}^{(1)},\ldots,\bar{\xv}^{(s)}$ be the optimal primal solution of the relaxation \eqref{eq:lagrelax} at this node. Then $\Gamma(\bar{\xv}^{(0)}) \ge \OPT$.
\end{lemma}

\begin{proof}
Let $\LAG(N)$ denote the exact optimal value of the relaxation
\eqref{eq:lagrelax} associated with node $N$ in the BP tree.
% By definition of an $\varepsilon$-approximate dual solution, we have
% $\LAG_\varepsilon(N) \in \LAG(N) + [0,\varepsilon].$
During the execution of BP, the value $\LAG(N)$ is used as the
node-selection criterion under the best-bound rule.

We claim that every internal node $N$ explored by BP satisfies
$\LAG(N) \ge \OPT.$
By contradiction, suppose there exists a node $N$ with
$\LAG(N) < \OPT$ that is selected for branching. By the best-bound
rule, this implies that for every other open node $N'$ in the BP tree, $\LAG(N') \le \LAG(N) < \OPT$.
On the other hand, there exists at least one node $N^*$ (open or pruned) whose
relaxation contains the globally optimal integral solution of the original
problem. For such a node, we have
$\OPT \le \LAG(N^*).$
Therefore, $N^*$ cannot be an open node at the time $N$ is selected, since all
open nodes have $\LAG(\cdot) < \OPT$. Hence, $N^*$ must already be
pruned.
If $N^*$ was pruned by integrality, then an incumbent solution of value $\OPT$
has already been found. If $N^*$ was pruned by bounds, then an incumbent of
value at least $\OPT$ must have been identified prior to pruning. In either
case, the current incumbent value is $\OPT$.
However, since $\LAG(N) < \OPT$, node $N$ would be pruned by the
bounding rule and thus could not be selected for branching, contradicting our
assumption. We conclude that every internal node of the BP tree satisfies
$\LAG(N) \ge \OPT$.

To conclude the proof, we show that for any node $N$ in the BP tree, the optimal
primal solution
$\bar{\xv}^{(0)}, \bar{\xv}^{(1)}, \ldots, \bar{\xv}^{(s)}$
of its relaxation~\eqref{eq:lagrelax} has $\Gamma$-value at least $\LAG(N)$, namely,
\begin{align}
\Gamma(\bar{\xv}^{(0)})
\;\ge\;
s \cdot \iprod{\cv^{(0)}}{\bar{\xv}^{(0)}}
+ \sum_{r \in [s]} \iprod{\cv^{(r)}}{\bar{\xv}^{(r)}} .
\label{eq:gammaBB}
\end{align}
To prove this, we expand the definition of $\Gamma$ and use the
fact that $\xv_*^{(r)}$ maximizes
$\iprod{\cv^{(r)} - (A^{(r)})^{\top} \mu_*^{(r)}}{\xv^{(r)}}$ to obtain
\begin{align}
\Gamma(\bar{\xv}^{(0)})
&\ge
s \cdot \iprod{\cv^{(0)} - (A^{(0)})^{\top} \overline{\mu}_*}{\bar{\xv}^{(0)}}
+ \sum_{r \in [s]} \iprod{\cv^{(r)} - (A^{(r)})^{\top} \mu_*^{(r)}}{\bar{\xv}^{(r)}}
+ \sum_{r \in [s]} \iprod{\mu_*^{(r)}}{\bv^{(r)}} \notag\\[4pt]
&=
s \cdot \iprod{\cv^{(0)}}{\bar{\xv}^{(0)}}
+ \sum_{r \in [s]} \iprod{\cv^{(r)}}{\bar{\xv}^{(r)}}
- \sum_{r \in [s]} \iprod{\mu_*^{(r)}}{A^{(0)} \bar{\xv}^{(0)} + A^{(r)} \bar{\xv}^{(r)}}
+ \sum_{r \in [s]} \iprod{\mu_*^{(r)}}{\bv^{(r)}}. \label{eq:gammaNode}
\end{align}
Since the solution
$\bar{\xv}^{(0)}, \bar{\xv}^{(1)}, \ldots, \bar{\xv}^{(s)}$
is feasible for the original constraints, i.e.,
\[
A^{(0)} \bar{\xv}^{(0)} + A^{(r)} \bar{\xv}^{(r)} \le \bv^{(r)}
\quad \text{for all } r \in [s],
\]
and $\mu_*^{(r)} \ge 0$, the last two terms of \eqref{eq:gammaNode} sum to a nonnegative value. Therefore,
\[
\Gamma(\bar{\xv}^{(0)})
\ge
s \cdot \iprod{\cv^{(0)}}{\bar{\xv}^{(0)}}
+ \sum_{r \in [s]} \iprod{\cv^{(r)}}{\bar{\xv}^{(r)}},
\]
which proves~\eqref{eq:gammaBB}. This completes the proof of the lemma.
\end{proof}

\end{document}
